\documentclass[12pt, letterpaper]{hmcpset}
\usepackage{math-19-preamble}

\name{Jayden Thadani}
\class{Calculus}
\assignment{Calculus on Manifolds, Problems 1 --- Functions on Euclidean Space}

\begin{document}

\begin{problem}[1]
	Prove that
	\[
		\norm{\bvec x} \le \sum_{i = 1}^n \abs{x_i}
	\]
	for $\bvec x = (x_1, \ldots, x_n) \in \bb R^n$.
\end{problem}

\begin{solution}
	We can prove this by induction on $n$ for $\bvec x = (x_{1}, \dots, x_{n}) \in \bb R^n$ \\\\
	For $n = 1$, we have
	\[
		\norm{\bvec x} = \abs{x_1} \le \abs{x_1}.
	\]\\
	Suppose for all $\bvec x \in \bb R^k$ we have
	\[
		\norm{\bvec x} \le \sum_{i = 1}^k \abs{x_i}.
	\]
	Then, for any $\bvec x \in \bb R^{k + 1}$ we can say that
	\[
		\begin{split}
			\norm{\bvec x}^2 & = \norm{(x_{1}, \dots, x_{k})}^2 + x_{k + 1}^2 \\
						 & \le \Big ( \sum_{i = 1}^k \abs{x_i} \Big )^2 + x_{k + 1}^2 \\
						 & \le \Big ( \sum_{i = 1}^k \abs{x_i} \Big )^2 + \abs{x_{k + 1}}^2 + 2 \abs{x_{k + 1}} \sum_{i = 1}^k \abs{x_i} \\
						 & = \Big ( \abs{x_{k + 1}} + \sum_{i = 1}^k \abs{x_i} \Big )^2.
		\end{split}
	\]\\
	That is,
	\[
		\norm{\bvec x}^2 \le \Big ( \sum_{i = 1}^{k + 1} \abs{x_i} \Big )^2
	\]
	and it follows for all $\bvec x \in \bb R^n$ for any $n \in \bb N$ that
	\[
		\norm{\bvec x} \le \sum_{i = 1}^n \abs{x_i}.
	\]
\end{solution}

\pagebreak

\begin{problem}[10]
	If $T : \bb R^n \to \bb R^m$ is a linear transformation, show that there is a number $M$ such that $\norm{T \bvec v} \le M \norm{\bvec v}$ for all $\bvec v \in \bb R^n$.
\end{problem}

\begin{solution}
	For the standard basis (or any orthonormal basis) $\bvec e_1, \ldots, \bvec e_n$ let us have $M_1, \ldots, M_n \in \bb R$ such that
	\[
		\norm{T \bvec e_i} = M_i \norm{\bvec e_i} = M_i.
	\]
	Then for any $\bvec v \in \bb R^n$ we have $\bvec v = \alpha_1 \bvec e_1 + \cdots + \alpha_n \bvec e_n$. This gives us
	\[
		\norm{\bvec v}^2 = \alpha_1^2 + \cdots + \alpha_n^2
	\]
	and
	\[
		\begin{split}
			\norm{T \bvec v} & = \norm{\alpha_1 T \bvec e_1 + \cdots + \alpha_n T \bvec e_n} \\
					   & \le \abs{\alpha_1} \norm{T \bvec e_1} + \cdots + \abs{\alpha_n} \norm{T \bvec e_n} \\
					   & = \abs{\alpha_1} M_1 + \cdots \abs{\alpha_n} M_n.
		\end{split}
	\]
	Let $m = \max \set{M_1, \ldots, M_n}$. Then we can say
	\[
		\begin{split}
			\norm{T \bvec v} & \le (m, \ldots, m) \cdot (\abs{\alpha_1}, \ldots, \abs{\alpha_n}) \\
					 & \le \norm{(m, \ldots, m)} \norm{\abs{\alpha_1}, \ldots, \abs{\alpha_n}} \\
					 & = \abs{m} \sqrt{n} \norm{\bvec v}.
		\end{split}
	\] \\
	Thus, $M = \abs m \sqrt n$ is the desired number that bounds $\norm{T \bvec v}$.
\end{solution}

\pagebreak

\begin{problem}[12]
	Let $\bb (R^n)^\vee$ denote the dual space of $\bb R^n$. If $\bvec x \in \bb R^n$, define $\varphi_{\bvec x} \in (\bb R^n)^\vee$ by $\varphi_{\bvec x}(\bvec y) = \bvec x \cdot \bvec y$. Define $T : \bb R^n \to (\bb R^n)^\vee$ by $T \bvec x = \varphi_{\bvec x}$. \\\\
	Show that $T$ is an invertible linear transformation and conclude that every $\varphi \in (\bb R^n)^\vee$ is $\varphi_{\bvec x}$ for a unique $\bvec x \in \bb R^n$.
\end{problem}

\begin{solution}
	If $T \bvec x = \bvec 0$ we must have $\varphi_{\bvec x} (\bvec y) = 0$ for all $\bvec y \in \bb R^n$. Thus, $\varphi_{\bvec x}(\bvec x) = 0$, giving us $\bvec x \cdot \bvec x = 0$ and thus, $\bvec x = \bvec 0$. Thus, $T$ is injective. \\\\
	Since $\bb R^n$ and $(\bb R^n)^\vee$ are both $n$-dimensional, $T$ is also invertible. Thus, for each $\varphi \in (\bb R^n)^\vee$ we have a unique $\bvec x \in \bb R^n$ such that $T \bvec x = \varphi$, and $\varphi = \varphi_{\bvec x}$.
\end{solution}

\pagebreak

\begin{problem}[14]
	Prove that the union of any (even infinite) number of open sets is open. Prove that the intersection of two (and hence, finitely many) open sets is open. Give a counter example for infinitely many open sets.
\end{problem}

\begin{solution}
	Suppose we have a collection of open sets, $\set{U_{\alpha}}$, indexed by $\alpha \in A$. For each $\bvec x \in \bigcup_{\alpha \in A} U_\alpha$ we have $\bvec x \in U_{\beta}$ for some $\beta \in A$. Since $U_\beta$ is open, we have $\bvec x \in I \subset U_\beta$ where $I$ is an open rectangle, and since we are looking at the union, we have $\bvec x \in I \subset \bigcup_{\alpha \in A} U_\alpha$. \\\\
	Suppose we have the open sets $U_1$ and $U_2$. For any $\bvec x \in U_1 \cap U_2$ we must have $\bvec x \in I_1 \subset U_1$ and $\bvec x \in I_2 \subset U_2$ where $I_1$ and $I_2$ are open rectangles. This gives us $\bvec x \in I_1 \cap I_2 \subset U_1 \cap U_2$. \\\\
	Note that the intersection of any two non-disjoint rectangles is also an open rectangle --- if we have rectangles $I_1 = \prod_{i = 1}^n (a_i, b_i)$ and $I_2 = \prod_{i = 1}^n (c_i, d_i)$ that have $\bvec x \in I_1 \cap I_2$, then we have $a_i, c_i < x_i < b_i, d_i$ for all $i \in \bb N_n$. Thus, we have
	\[
		I_1 \cap I_2 = \prod_{i = 1}^n (\max \set{a_i, c_i}, \min \set{b_i, d_i})
	\]
	as an open rectangle. \\\\
	Since this is an open rectangle, we have the desired open rectangle subset containing any point in the intersection of $U_1$ and $U_2$. By induction we can extend this to any finite intersection of open sets. \\\\
	This result doesn't hold for infinitely many open sets --- the intersection of all open intervals of the form $(-1/n, 1/n)$ is $\set{0}$, which is not open.
\end{solution}

\pagebreak

\begin{problem}[15]
	Prove that $\set{\bvec x \in \bb R^n : \norm{\bvec x - \bvec a} < r}$ is open (for any $\bvec a \in \bb R^n$ and any $r > 0$).
\end{problem}

\begin{solution}
	Let us call the set $N_r(\bvec a)$. For any $\bvec x \in N_r(\bvec a)$, consider $\rho = r - \norm{\bvec x - \bvec a}$. Then by the triangle inequality we have $N_\rho(\bvec x) \subset N_r(\bvec a)$. By geometry (the Pythagorean theorem), we have
	\[
		\prod_{i = 1}^n \Big ( x_i - \frac{\rho}{\sqrt 2}, x_i + \frac{\rho}{\sqrt 2} \Big ) \subset N_\rho(\bvec x) \subset N_r(\bvec a).
	\]
	This is the open rectangle subset (obviously containing $\bvec x$) required for $N_r(\bvec a)$ to be open.
\end{solution}

\pagebreak

\begin{problem}[19]
	If $A$ is a closed set that contains every rational number $r \in [0, 1]$, then show that $[0, 1] \subset A$.
\end{problem}

\begin{solution}
	Note that for $A$ to be closed, we must have $\bb R \setminus A$ open. Suppose, by way of contradiction, that we didn't have $[0, 1] \subset A$, and instead, $x \in [0, 1]$ but $x \notin A$. \\\\
	Since there is a rational number beween any two real numbers, any open interval $(a, b)$ with $x \in (a, b)$ would contain a rational number in $[0, 1]$. Thus, there would be no open interval $(a, b)$ with $x \in (a, b) \subset \bb R \setminus A$, making it not open, and giving rise to a contradiction. Thus, our supposition that there is such an $x$ must be false.
\end{solution}

\pagebreak

\begin{problem}[21]
	\begin{enumerate}
		\item If $A \subset \bb R^n$ is closed and $\bvec x \notin A$ prove that there is a number $d > 0$ such that $\norm{\bvec y - \bvec x} \ge d$ for all $\bvec y \in A$.
		\item If, for $A, B \subset \bb R^n$, $A$ is closed, $B$ is compact, and $A \cap B = \emptyset$, prove that there is $d > 0$ such that $\norm{\bvec y - \bvec x} \ge d$ for all $\bvec y \in A$ and $\bvec x \in B$.
		\item Give a counterexample to part (b) in $\bb R^2$ for if $A$ and $B$ are both closed but neither is compact.
	\end{enumerate}
\end{problem}

\begin{solution}
	\begin{enumerate}
		\item Since $\bvec x \in \bb R^n \setminus A$, an open set, we have an open rectangle $I$ such that $\bvec x \in I \subset \bb R^n \setminus A$. For $I = \prod_{i = 1}^n (a_i, b_i)$, define
			\[
				d = \min \set{\abs{x - a_i}, \abs{x - b_i} : i \in \bb N_n}.
			\]
			This allows to contain the open ball $N_d(\bvec x) = \set{\bvec y \in \bb R^n : \norm{\bvec y - \bvec x} < d}$ inside $I$. That is, the coordinates of any point in $N_d(\bvec y)$ are displaced by less than the minimum required to leave $I$, and thus, $N_d(\bvec x) \subset I$.\\\\
			Since $N_d(\bvec x) \subset I \subset \bb R^n \setminus A$, we have $\bvec y \notin A$ for any point $\bvec y$ with $\norm{\bvec y - \bvec x} < d$, and the contrapositive is the desired property for $d$.
		\item Since $A$ is closed, for any $\bvec x \in B$, there exists some $d'$ that gives us that $N_{d'} \cap A$ is empty. Consider $\mathcal O$, the open cover of $B$ given by considering a $d'/2$-neighbourhood of each point. By the compactness of $B$, we can reduce this to a finite sub-cover $\bar{\mathcal O} = \set{U_1, \ldots, U_n}$ where each $U_i$ is a $d'_i/2$-neighbourhood of some point in $B$.\\\\
			For any point $\bvec x \in B$, we have $\bvec x \in U_i \in \bar{\mathcal O}$, and thus, $\bvec x$ is at distance greater than $d_i'/2$ from any $\bvec y \in A$. Since there are only finitely many $U_i$, there are only finitely many corresponding $d_i'$ and thus, we can bound $d$.
			\[
				d > \min \set{d_i' : i \in \bb N_n} > 0.
			\]
			That is, we do not have arbitrarily close points of $A$ and $B$.
		\item Consider $A = \set{(x, y) \in \bb R^2 : y \ge 1/x, x > 0}$ and $B = \set{(x, y) \in \bb R^2 : y \le - 1/x, x > 0}$. Since $1/0$ is undefined, we do not consider any points for which $x = 0$ to be part of either set. \\\\
			Note that both sets are closed. For any $(x, y) \in \bb R^2 \setminus A$ there is a known technique to find the point on the curve $y = 1/x$ nearest to $(x, y)$. The neighbourhood $N_s(x, y)$ where $s$ is the distance of this nearest point is then contained in $\bb R^2 \setminus A$. Since neighbourhoods are open, this makes $\bb R^2 \setminus A$ open, and thus, $A$ closed. Similarly, we can show that $B$ is closed. Neither of them is compact since neither of them is bounded. \\\\
			Note finally, that for any $d > 0$, there is some $2/n < d$ where $n \in \bb N$. Thus, for any $d > 0$ we have the points $(n, 1/n)$ and $(n, -1/n)$ at distance $2/n < d$.
	\end{enumerate}
\end{solution}

\pagebreak

\begin{problem}[22]
	If (all subsets of $\bb R^n$), $U$ is open, and $C \subset U$ is compact, then show that there is a compact set $D$ such that $C \subset D^\circ$ and $D \subset U$ ($D^\circ$ denotes the interior of $D$).
\end{problem}

\begin{solution}
	Consider $\mathcal O$, the open cover of $C$ produced by looking at the open rectangles $I \subset U$ containing each point in $C$. Let $\bar{\mathcal O}$ be the finite subcover that $\mathcal O$ reduces to. Note that the union of the finite subcover is an open subset of $U$ containing $C$. To be exact,
	\[
		C \subset \bigcup_{V \in \bar{\mathcal O}} V \subset U.
	\]
	More over, since it is the union of finitely many open sets, it is a bounded set. Thus, its closure is a closed and bounded, and therefore compact subset $D \subset U$ and it is $D^\circ$, the open interior that contains $C$.
\end{solution}

\pagebreak

\begin{problem}[23]
	If $f : A \to \bb R^m$ and $\bvec a \in A$ show that $\lim_{\bvec x \to \bvec a} f(\bvec x) = \bvec b$ if and only if $\lim_{\bvec x \to \bvec a} f_i(\bvec x) = b_i$ for $i \in \bb N_m$. Here, $f = (f_1, \ldots, f_m)$ and $\bvec b = (b_1, \ldots, b_m)$.
\end{problem}

\begin{solution}
	Suppose we have $\lim_{\bvec x \to \bvec a} f_j(\bvec x) \neq b_j$ for some $j$. Then for some $\varepsilon > 0$ no matter what $\delta > 0$ we choose, there will be some point $\bvec x$ with $\norm{\bvec x - \bvec a} < \delta$ and ${\abs{f_j(\bvec x) - f_j(\bvec a)} \ge \varepsilon}$. Thus, for that same $\varepsilon$, no matter, what $\delta > 0$ we choose, we will have, at the same point $\bvec x$, the following.
	\[
		\begin{split}
			\norm{f(\bvec x) - f(\bvec a)}^2 & = (f_1(\bvec x) - f_1(\bvec a))^2 + \cdots + (f_m(\bvec x) - f_m(\bvec a))^2 \\
							 & \ge (f_j(\bvec x) - f_j(\bvec a))^2 \\
							 & \ge \varepsilon^2
		\end{split}
	\]
	which means
	\[
		\norm{f(\bvec x) - f(\bvec a)} \ge \varepsilon.
	\]\\
	That is, for some $\varepsilon > 0$, no matter what $\delta > 0$ we choose there will be some point $\bvec x$ with $\norm{\bvec x - \bvec a} < \delta$ and $\norm{f(\bvec x) - f(\bvec a)} \ge \varepsilon$. Thus, $\lim_{\bvec x \to \bvec a} f(\bvec x) \neq \bvec b$.\\\\
	Suppose we have $\lim_{\bvec x \to \bvec a} f_i(\bvec x) = b_i$ for all $i \in \bb N_m$ and $\varepsilon > 0$. Then let $\varepsilon' = \frac{\varepsilon}{\sqrt m}$.\\\\
	For each $i \in \bb N_m$ we have some $\delta_i$ such that $\abs{f_i(\bvec x) - f_i(\bvec a)} < \varepsilon'$ for all $\bvec x$ with $\norm{\bvec x - \bvec a} < \delta_i$. Let $ \delta' = \min \set{\delta_i}$. Then for all $\bvec x$ with $\norm{\bvec x - \bvec a} < \delta'$ we have
	\[
		\begin{split}
			\norm{f(\bvec x) - f(\bvec a)}^2 & = (f_1(\bvec x) - f_1(\bvec a))^2 + \cdots + (f_m(\bvec x) - f_m(\bvec a))^2 \\
						       & < \varepsilon'^2 + \cdots + \varepsilon'^2 \quad \text{($m$ times)} \\
						       & = m \frac{\varepsilon^2}{m} \\
						       & = \varepsilon^2.
		\end{split}
	\]
	and thus,
	\[
		\norm{f(\bvec x) - f(\bvec a)} < \varepsilon.
	\]
	That is,
	\[
		\lim_{\bvec x \to \bvec a} f(\bvec x) = \bvec b.
	\]
\end{solution}

\pagebreak

\begin{problem}[24]
	Prove that $f : A \to \bb R^m$ is continuous at $\bvec a \in A$ if and only if each $f_i$ is. Again, $f = (f_1, \ldots, f_m)$.
\end{problem}

\begin{solution}
	Let $\bvec b = (b_1, \ldots, b_m) = f(\bvec a)$. Then $f_i(\bvec a) = b_i$ for all $i \in \bb N_m$.\\\\
	We know that $\lim_{\bvec x \to \bvec a} f(\bvec x) = \bvec b$ if and only if $\lim_{\bvec x \to \bvec a} f_i(\bvec x) = b_i$ for all $i \in \bb N_m$. Thus, $f$ is continuous if and only if each $f_i$ is.
\end{solution}

\pagebreak

\begin{problem}[25]
	Prove that a linear transformation $T : \bb R^n \to \bb R^m$ is continuous.
\end{problem}

\begin{solution}
	Consider any point $T \bvec a \in \bb R^m$ and any $\varepsilon > 0$. We can find points $\bvec x \in \bb R^n$ such that $\norm{T \bvec x - T \bvec a} < \varepsilon$ since we have bounded $T \bvec v$ by $\norm{\bvec v}$. \\\\
	Specifically, if we have $\norm{T \bvec v} \le M \norm{\bvec v}$, then we can choose $\delta = \varepsilon/M$. If we look at $\bvec x$ such that $\bvec v = \bvec x - \bvec a$ has $\norm{\bvec v} < \delta$, then we have
	\[
		\begin{split}
			\norm{T \bvec x - T \bvec a} & = \norm{T \bvec v} \\
						     & \le M \norm{\bvec v} \\
						     & < M \frac{\varepsilon}{M} \\
						     & = \varepsilon
		\end{split}
	\]\\
	That is, for all $\bvec x \in \bb R^n$ with $\norm{\bvec x - \bvec a} < \delta$ we have $\norm{T \bvec x - T \bvec a} < \varepsilon$. $T$ is continuous everywhere.
\end{solution}

\pagebreak

\begin{problem}[26]
	Let $A = \set{(x, y) \in \bb R^2 : x > 0 \text{ and } 0 < y < x^2 }$.
	\begin{enumerate}
		\item Show that every straight line through $\bvec 0$ contains an interval around $\bvec 0$ which is in $\bb R^2 \setminus A$.
		\item Define $f : \bb R^2 \to \bb R$ by
			\[
				f(\bvec v) = \begin{cases}
					0 & \bvec v \notin A \\
					1 & \bvec v \in A.
				\end{cases}
			\]
			For $\bvec v \in \bb R^2$ define $g_{\bvec v} : \bb R \to \bb R$ by $g_{\bvec v}(\lambda) = f(\lambda \bvec v)$. Show that each $g_{\bvec v}$ is continuous at $0$, but $f$ is not continuous at $\bvec 0$.
	\end{enumerate}
\end{problem}

\begin{solution}
	\begin{enumerate}
		\item Every straight line through $\bvec 0$ is of the form $y = \alpha x$ for some $\alpha \in \bb R$, or $x = 0$.\\\\
			For $x = 0$, the whole line is not in $A$, obviously. \\\\
			For all other lines consider points in the graph of the line with $x \in (-\abs \alpha, \abs \alpha)$. These points give us $x^2 < \alpha x$ and thus, $x^2 < y$ for positive $x$ and positive $\alpha$. For positive $x$ and negative $\alpha$ we have $y < 0$ and nonpositive $x$ are already excluded from $A$. Thus, for $x$ in this interval, we have points in an interval about $\bvec 0$ that are in $\bb R^2 \setminus A$.
		\item Choose $\delta = \frac{v_2}{v_1}$. Then, for any $\varepsilon > 0$, for all $\lambda$ with $\abs{\lambda - 0} < \delta$ we have $\abs{g_{\bvec v}(\lambda) - g_{\bvec v}(0)} = 0 < \varepsilon$. Thus, $g_{\bvec v}$ is continuous at $0$.\\\\
			However, for any $0 < \delta$, there are always points $(x, y) = (\frac{\delta}{\sqrt 2}, \frac{\delta}{2 \sqrt 2})$ with $\norm{(x, y) - \bvec 0} = \delta \sqrt{\frac{5}{8}} < \delta$ and $x^2 = \frac{\delta}{2} > \frac{\delta}{2 \sqrt 2} = y$. These points $(x, y)$ are in $A$ and give $f(x,y) = 1$.\\\\
			Thus, for any $0 < \varepsilon < 1$ there is no $\delta > 0$ such that for all $\bvec v$ with $\norm{\bvec v - \bvec 0} < \delta$ we have $\abs{f(\bvec v) - f(\bvec 0)} < \varepsilon$. That is, $f$ is not continuous at $\bvec 0$.
	\end{enumerate}
\end{solution}

\pagebreak

\begin{problem}[27]
	Prove that $\set{\bvec x \in \bb R^n : \norm{\bvec x - \bvec a} < r}$ is open by considering the function $f : \bb R^n \to \bb R$ defined by $f(\bvec x) = \norm{\bvec x - \bvec a}$ (for $\bvec a \in \bb R^n$, $r > 0$).
\end{problem}

\begin{solution}
	First we will show that $f$ is continuous at all $\bvec b \in \bb R^n$.\\\\
	For any $\varepsilon > 0$, if we want $\abs{f(\bvec b) - f(\bvec x)} < \varepsilon$, then we are looking for points with $\abs{\norm{\bvec b - \bvec a} - \norm{\bvec x - \bvec a}} < \varepsilon$. The triangle inequality tells us that if we have $\norm{\bvec x - \bvec b} < \varepsilon$ then we will always have $\norm{\bvec x - \bvec a} \le \norm{\bvec x - \bvec b} + \norm{\bvec b - \bvec a}$ and thus, $\norm{\bvec x - \bvec a} - \norm{\bvec b - \bvec a} \le \norm{\bvec x - \bvec b} < \varepsilon$. The triangle inequality also gives us $\norm{\bvec b - \bvec a} \le \norm{\bvec x - \bvec a} + \norm{\bvec x - \bvec b}$ and thus, $\norm{\bvec b - \bvec a} - \norm{\bvec x - \bvec a} \le \norm{\bvec x - \bvec b} < \varepsilon$. \\\\
	Thus, choosing $\delta = \varepsilon$ gives us that for every $\bvec x$ with $\norm{\bvec x - \bvec b} < \delta$ we have $\abs{f(\bvec x) - f(\bvec b)} < \varepsilon$. That is, $f$ is continuous.\\\\
	Consider the pre-image of the open interval $(-r, r)$ under $f$. That is, consider $\set{\bvec x \in \bb R^n : \norm{\bvec x - \bvec a} <r}$. Since $f$ is continuous and $(-r, r)$ is open, $f^\text{pre}((-r, r))$ is open, as desired.
\end{solution}

\pagebreak

\begin{problem}[29]
	If $A \subset \bb R^n$ is compact, prove that every continuous function $f : A \to \bb R$ takes on a maximum and minimum value.
\end{problem}

\begin{solution}
	Note that we have $f^{\text{img}}(A)$ as compact. Thus it is closed and bounded, and we can talk about $\sup f^{\text{img}}(A)$ and $\inf f^{\text{img}}(A)$. We claim that these are contained in $f^{\text{img}}(A)$ and are thus its respective maximum and minimum values.\\\\
	Since $f^{\text{img}}(A)$ is closed, it's complement must be open. If $\sup f^\text{img}(A)$ were not in $A$, then it would be in the complement and thus, the complement would contain a point around which it does not contain any open interval making it not open --- any open interval containing $\sup f^\text{img}(A)$ would contain points less than $\sup f^\text{img}(A)$, and thus, would contain points of $A$. Thus, $\sup f^\text{img}(A)$ must be in $A$. Similarly, $\inf f^\text{img}(A)$ must be in $A$. \\\\
	Then, by their definition they are the maximum and minimum values that $f$ takes on $A$.
\end{solution}

\end{document}
